%% ============================================================
%% nro/avaliacao.tex — o barema do professor e as competências
%%
%% Mantido pela DIRETORIA, com o professor responsável. O barema é o
%% mesmo da autoavaliação do membro (Parte IV) e do parecer do professor
%% (Parte V): quem escreve sabe, desde o primeiro rascunho, como vai ser
%% avaliado. Os pesos somam 100.
%% ============================================================

%% \criterio{chave}{peso}{nome}{o que se espera para a nota máxima}
\criterio{aderencia}{15}{Aderência à formação}
  {As atividades são de engenharia do curso do estudante, e a relação com as disciplinas e com as competências está demonstrada, não apenas afirmada.}
\criterio{profundidade}{25}{Profundidade técnica}
  {O desenvolvimento mostra método, cálculos, decisões justificadas e domínio das ferramentas, num nível comparável ao de um relatório de iniciação científica.}
\criterio{contribuicao}{20}{Contribuição individual}
  {Fica claro o que o estudante fez, decidiu e entregou, separado do trabalho da equipe, e os registros e os supervisores confirmam.}
\criterio{evidencias}{15}{Evidências e rastreabilidade}
  {Cada atividade tem evidências verificáveis (códigos no SOMA, repositórios, arquivos de projeto, documentos emitidos), e as horas são consistentes com elas.}
\criterio{resultados}{15}{Resultados e impacto}
  {Há resultados concretos, medidos e comparados com o esperado, e o impacto no projeto, na equipe ou na comunidade está demonstrado.}
\criterio{comunicacao}{10}{Comunicação}
  {O texto é claro, organizado e correto; figuras e tabelas têm legenda e fonte; a linguagem técnica é precisa.}

%% As competências gerais do egresso de Engenharia, na ordem dos incisos
%% do art. 4º da Resolução CNE/CES nº 2, de 24 de abril de 2019 (as DCN
%% dos cursos de Engenharia). A atividade cita as suas em
%% competencias = {I, III, V}.
\definircompetencia{I}{Formular e conceber soluções desejáveis de engenharia, analisando e compreendendo os usuários dessas soluções e o seu contexto.}
\definircompetencia{II}{Analisar e compreender os fenômenos físicos e químicos por meio de modelos simbólicos, físicos e outros, verificados e validados por experimentação.}
\definircompetencia{III}{Conceber, projetar e analisar sistemas, produtos (bens e serviços), componentes ou processos.}
\definircompetencia{IV}{Implantar, supervisionar e controlar as soluções de engenharia.}
\definircompetencia{V}{Comunicar-se eficazmente nas formas escrita, oral e gráfica.}
\definircompetencia{VI}{Trabalhar e liderar equipes multidisciplinares.}
\definircompetencia{VII}{Conhecer e aplicar com ética a legislação e os atos normativos no âmbito do exercício da profissão.}
\definircompetencia{VIII}{Aprender de forma autônoma e lidar com situações e contextos complexos, atualizando-se em relação aos avanços da ciência, da tecnologia e aos desafios da inovação.}
